% ==============================================================================
% ПАРТИЯ B03-036: DISS-005-B1958 .. DISS-005-B2011
% Глава 3: Методология и автоматизация аннотирования осциллограмм
% ==============================================================================

% DISS-005-B1958
\section{Разметка (аннотирование) осциллограмм}\label{sec:ch3_oscillogram_annotation}

% DISS-005-B1959
После того как осциллограммы собраны, унифицированы, деперсонализированы и нормализованы, следующим логическим шагом является их разметка или аннотирование. Этот процесс заключается в присвоении метки каждому участку или всей осциллограмме, характеризующей тип произошедшего события (или его отсутствие). Качественная разметка является фундаментом для последующего глубокого анализа данных, обучения моделей машинного обучения и решения специфических задач релейной защиты и автоматики.

% DISS-005-B1960
{\emergencystretch=1.5em Для диссертационного исследования, сфокусированного на устройствах быстродействующего автоматического ввода резерва (БАВР), разметка осциллограмм играет ключевую роль. Хотя можно моделировать логику работы БАВР и проверять ее на синтетических данных, реальные осциллограммы предоставляют бесценную информацию о поведении энергосистемы и самого устройства БАВР в реальных условиях эксплуатации. Безусловно, первичную разметку можно провести за счёт моделирования работы современных методов работы органов РЗА, а в случае нашего исследования~-- современных органов реле направления мощности (РНМ). Для первых этапов это и будет осуществлять, что полезно для обучения и последующего сравнения различных методов РНМ.\par}

% DISS-005-B1961
Но это не всегда позволяет учесть все нюансы реального поведения оборудования и энергосистемы, особенно в сложных или нештатных ситуациях (Глава 1х\commentnote{DISS-005-C0130}{Алексей Евдаков [2]}{Как будет подготовлен и нормально описан первый раздел, надо стедлать ссылку}{х}). Поэтому качественная разметка очень важна, и она позволяет:

\begin{itemize}
% DISS-005-B1962
  \item \textbf{Идентифицировать и классифицировать события.} То есть выделить из общего массива данных те осциллограммы, которые содержат интересующие аварийные режимы (например, короткие замыкания различных видов), аномальные состояния (например, однофазные замыкания на землю, просадки напряжения) или оперативные переключения, релевантные для анализа работы БАВР.
% DISS-005-B1963
  \item \textbf{Формировать наборы данных для верификации алгоритмов БАВР.} Создавать специализированные выборки осциллограмм, отражающие сценарии, при которых БАВР должен или не должен срабатывать. Это особенно важно для оценки корректности функционирования измерительных органов, таких как реле направления мощности, и логики принятия решений устройством.
% DISS-005-B1964
  \item \textbf{Проводить статистический анализ реальных явлений.} Собранный и размеченный датасет открывает возможности для исследования характеристик реальных переходных процессов. Например, как будет подробно описано в последующих главах (Глава \commentnote{DISS-005-C0131}{Алексей Евдаков [2]}{Добавить ссылку на исследование ОЗЗ}{\commentnote{DISS-005-C0132}{Алексей Евдаков [2]}{Ну и пока ещё в планах}{ХХ}}), на основе данного датасета был проведен анализ перенапряжений при однофазных замыканиях на землю (ОЗЗ) в реальных сетях с применением классических датчиков.
% DISS-005-B1965
  \item \textbf{Создавать основу для обучения ML-моделей.} Размеченные данные необходимы для обучения моделей машинного обучения, способных автоматизировать классификацию событий, предсказывать поведение системы или даже участвовать в логике работы адаптивных устройств РЗА.
\end{itemize}

% DISS-005-B1966
Важно отметить, что создаваемый датасет, благодаря своей обширности и разнообразию, может быть использован не только для исследований, связанных с БАВР. Он представляет ценность для анализа широкого круга явлений в энергосистемах, таких как оценка бросков тока при включении силовых трансформаторов, анализ пусковых режимов двигателей, исследование влияния различных видов нагрузок на качество электроэнергии и~т.\,д. При необходимости, для конкретных осциллограмм из датасета можно запросить дополнительную информацию о параметрах оборудования из частных источников (с соблюдением требований конфиденциальности).

% DISS-005-B1967
\subsection{Необходимость разметки для обучения ML моделей и анализа}\label{sec:ch3_annotation_necessity}

% DISS-005-B1968
Ручная разметка осциллограмм~-- это трудоемкий и затратный по времени процесс, требующий высокой квалификации и экспертизы от специалистов [Х\commentnote{DISS-005-C0133}{Алексей Евдаков [2]}{Ссылка на нашу статью про датасет. Но можно и кучу других приплести, очень простой и очевидный фактор всё же}{Х}]. Субъективность оценки и сложность точного определения границ событий являются неотъемлемыми вызовами этого подхода.

% DISS-005-B1969
{\emergencystretch=1.5em Для систематизации процесса аннотирования была предложена иерархическая структура, представленная на Рисунок~3.10. Эта структура делит все события на четыре основные группы:\par}

\begin{itemize}
% DISS-005-B1970
  \item \textbf{Группа 0: «Нет события» (No Event)}~-- Сегменты, где энергосистема работает в нормальном режиме, без значимых отклонений или событий. Сюда же относятся осциллограммы, содержащие только фоновый шум без полезного сигнала.
% DISS-005-B1971
  \item \textbf{Группа 1: «Оперативные переключения» (Operational Switching)}~-- Сегменты, отражающие штатные коммутации нагрузки, обычные операции с оборудованием (например, пуск двигателя, переключение обмоток трансформатора) или действия РЗА во время ликвидации аварии (например, работа выключателя).
% DISS-005-B1972
  \item \textbf{Группа 2: «Анормальные события» (Abnormal Events)}~-- События, которые отклоняются от нормального режима работы, но не требуют немедленного отключения (например, ОЗЗ в сетях с изолированной или компенсированной нейтралью, незначительные просадки напряжения).
% DISS-005-B1973
  \item \textbf{Группа 3: «Аварийные события» (Fault Events)}~-- События, которые вызывают или должны вызывать срабатывание устройств релейной защиты (например, короткие замыкания, перегрузки, ведущие к отключению) или в нашем случае~-- срабатывание БАВР.
\end{itemize}

\begin{figure}[htbp]
\centering
% DISS-005-TAB-086
% DISS-005-FIG-092
% DISS-005-B1974
\includegraphics[width=0.95\textwidth]{image95.png}
% DISS-005-B1975
\caption{Упрощенная схема иерархии классификации.}\label{fig:ch03_classification_hierarchy}
\end{figure}

% DISS-005-B1976
Отметим, что данная разметка носит общий характер для всеобъемлющего разбора данного набора данных. Более простая разметка для конкретных задач и органов РЗА будет описана в соответствующих главах.

% DISS-005-B1977
Пример декодирования метки ML-сигнала приведен на Рисунок~3.11.

\begin{figure}[htbp]
\centering
% DISS-005-TAB-087
% DISS-005-FIG-093
% DISS-005-B1978
\includegraphics[width=0.6\textwidth]{image96.png}
% DISS-005-B1979
\caption{Пример декодирования сигнала ML..}\label{fig:ch03_ml_signal_decoding}
\end{figure}

% DISS-005-B1980

% DISS-005-B1981
\subsection{Подготовка данных для ручной разметки}\label{sec:ch3_manual_annotation_prep}

% DISS-005-B1982
Для первоначального обучения моделей и валидации подходов была выполнена ручная разметка подмножества осциллограмм. Было отобрано 40\commentnote{DISS-005-C0134}{Алексей Евдаков [2]}{Может меняться}{9} осциллограмм, записанные в сетях 50~Гц с частотой дискретизации 1600~Гц, которые содержали относительно высокую частоту событий.

% DISS-005-B1983
Для удобства ручной разметки и последующего обучения моделей ML, исходные COMTRADE файлы были преобразованы в единый CSV-файл. Этот процесс включал:

\begin{itemize}
% DISS-005-B1984
  \item Итерацию по всем осциллограммам.
% DISS-005-B1985
  \item Разделение сигналов по секциям шин (с добавлением суффикса «\_BusXX» к имени осциллограммы).
% DISS-005-B1986
  \item Поскольку многие осциллограммы содержат лишь короткий участок с интересующим событием, а остальная часть является предаварийным или послеаварийным режимом, для экономии ресурсов и фокусировки на событии производилось «вырезание» областей: сохранялся участок события и регион, равный 20 периодам до и после события.
% DISS-005-B1987
  \item Если в осциллограмме присутствовало несколько событий находящихся на расстоянии больше 40 периодов, они также разделялись (с добавлением суффикса «\_eventNXX»).
% DISS-005-B1988
  \item Для осциллограмм без событий сохранялся 20-периодный участок из середины записи.
\end{itemize}

% DISS-005-B1989
\subsection{Автоматизация процесса разметки и перспективы}\label{sec:ch3_annotation_automation}

% DISS-005-B1990
Учитывая огромный объем данных (более 50 тысяч осциллограмм в текущем датасете), ручная разметка всего массива не представляется возможной. Поэтому ключевой задачей является разработка и применение методов автоматической разметки. На основе первоначально размеченного вручную подмножества могут быть обучены модели машинного обучения (например, свёрточные нейронные сети (CNN) или многослойные персептроны (MLP) на основе спектральных данных) для предварительной классификации и сортировки оставшихся осциллограмм. Этот подход будет подробно рассмотрен в следующих главах (Глава \commentnote{DISS-005-C0135}{Алексей Евдаков [2]}{Первичные наброски уже есть, и будут дополняться, так что хорошо бы включить и в диссертацию}{ХХ}) 

% DISS-005-B1991
Этот процесс может быть итерационным:

\begin{itemize}
% DISS-005-B1992
  \item Обучение начальной модели на малом размеченном наборе.
% DISS-005-B1993
  \item Применение модели для автоматической разметки большего неразмеченного набора.
% DISS-005-B1994
  \item Экспертная проверка и коррекция автоматически присвоенных меток, с особым вниманием к сомнительным или сложным случаям.
% DISS-005-B1995
  \item Дообучение или переобучение модели на расширенном и уточненном наборе данных.
% DISS-005-B1996
  \item Повторение шагов 2--4 до достижения приемлемого качества автоматической разметки.
\end{itemize}

% DISS-005-B1997
Такой циклический подход позволяет постепенно улучшать как качество самой разметки, так и точность моделей машинного обучения

% DISS-005-B1998
\subsection{Фильтрация «пустых» осциллограмм}\label{sec:ch3_empty_oscillograms_filter}

% DISS-005-B1999

% DISS-005-B2000

% DISS-005-B2001

% DISS-005-B2002

% DISS-005-B2003
На начальном этапе формирования датасета, особенно при сборе данных с пусконаладочных работ, значительную часть составляют записи, где не зафиксировано явных аварийных или переходных процессов. Хотя осциллограммы, отражающие нормальный режим работы, важны для обучения моделей отличать его от аномального, их избыточное количество (особенно если это просто фоновый шум или незначительные флуктуации без событий, представляющих интерес для анализа) затрудняет поиск и обработку действительно информативных записей. Для первичного отсеивания таких «истинно пустых» или не представляющих интереса осциллограмм была разработана специальная функция (EmptyOscFilter). Результаты работы данной функции приведены в таблице~2. Основная логика ее работы заключается в следующем.

% DISS-005-B2004
\noindent\textbf{\underline{Чтение осциллограммы}}\par\vspace{0.1cm}

% DISS-005-B2005
Загружается «.cfg» и «.dat» файл осциллограммы. Определяется частота дискретизации и, на ее основе, размер окна для БПФ (обычно один период основной частоты).

% DISS-005-B2006
\noindent\textbf{\underline{Анализ аналоговых каналов}}\par\vspace{0.1cm}

% DISS-005-B2007
Для каждого аналогового канала, соответствующего критериям (например, по имени содержащего «I»' или «U», что указывает на ток или напряжение):

\begin{itemize}
% DISS-005-B2008
  \item \textbf{Опциональная нормализация:} Если активирована опция нормализации (use\_\allowbreak norm\_\allowbreak osc) и~предоставлен файл с~данными коэффициентами (norm\_\allowbreak coef.csv), предпринимается попытка нормализовать сигнал с~помощью. Если нормализация невозможна, анализ продолжается с~«сырыми» данными (и другими уставками).
% DISS-005-B2009
  \item \textbf{Проверка «чистоты» начального участка:} Для «сырых» сигналов (или если нормализация не применялась) анализируется начальный участок сигнала (например, первый период) на предмет доминирования первой гармоники над высшими. Если сигнал слишком «грязный» (много шума/гармоник), он может считаться неактивным для дальнейшего анализа порогов (уставок).
% DISS-005-B2010
  \item \textbf{Расчет амплитуды первой гармоники ($H_1$):} Для всего сигнала или его репрезентативных участков рассчитывается временной ряд амплитуд первой гармоники.
% DISS-005-B2011
  \item \textbf{Применение уставок:} Рассчитанные статистики (максимальное значение $H_1$, дельта $H_1$, стандартное отклонение $H_1$) сравниваются с заданными уставоками. Пороги могут быть абсолютными (для нормализованных сигналов) или относительными (для «сырых» сигналов, относительно начальной амплитуды $H_1$). В таблице~6 приведены значения используемы уставок. Здесь delta~-- расхождение между максимальным и минимальным значением; std\_\allowbreak dev~-- стандартное отклонение; max\_\allowbreak abs~-- максимальное значение (для отсеивания шумов).
\end{itemize}
